\subsection{Pluriclosed metrics under deformation}
\label{sec:pluriclosed-deformations}

Let $M$ be a compact smooth manifold of real dimension $2n$, $n\geq1$, and let $J_t$, $t\in U\subset\mathbb R^\ell$, be a smooth family of integrable complex structures, where $U$ is an open neighborhood of $0$. Write $\Omega_t^{p,q}$ for the smooth complex forms of type $(p,q)$ for $J_t$. A pluriclosed form belongs to the kernel of $\partial_t\bar\partial_t$. The corresponding Aeppli group is
\[
 H_A^{1,1}(M,J_t)
 =\frac{\ker(\partial_t\bar\partial_t:\Omega_t^{1,1}\to\Omega_t^{2,2})}
 {\partial_t\Omega_t^{0,1}+\bar\partial_t\Omega_t^{1,0}},
 \qquad h_A^{1,1}(M,J_t)=\dim_{\mathbb C}H_A^{1,1}(M,J_t).
\]
\index{Aeppli cohomology}
Constancy of this dimension gives a smooth projection onto the closedness constraint. Positivity then persists near a given strictly positive form.

\Needspace{8\baselineskip}
\begin{lemma}\label{lem:aeppli-smooth-projection}
The number $h_A^{1,1}(M,J_t)$ is finite. If it is constant near $0$, then, after shrinking $U$, the orthogonal projections
\[
 \mathcal P_t:\Omega_t^{1,1}\longrightarrow\ker(\partial_t\bar\partial_t)
\]
for any smooth family of auxiliary Hermitian metrics vary smoothly with $t$ on smooth forms and preserve reality.
\end{lemma}
\begin{proof}
Suppress $t$ initially, and use the auxiliary metric to form adjoints. Put $C=\partial\bar\partial$ and $E=\bar\partial\partial^*$. The nonnegative self-adjoint operator
\begin{equation}\label{eq:aeppli-elliptic-operator}
 \mathcal L=C^*C+CC^*+E^*E+EE^*
                  +\partial\partial^*+\bar\partial\bar\partial^*
\end{equation}
is elliptic of order four. Indeed, at a nonzero real covector $\xi$, the symbols of $\partial$ and $\bar\partial$ are exterior multiplication by its $(1,0)$ and $(0,1)$ parts, up to the common factor $\mathrm i$. These are orthogonal and have equal norm $r>0$. Split the exterior algebra according to whether each of these two directions is present. The four fourth-order terms act respectively on the four possibilities, with sum $r^4\operatorname{Id}$. The lower-order terms do not affect ellipticity.

The quadratic form of $\mathcal L$ shows that
\[
 \ker\mathcal L=\ker C\cap\ker\partial^*\cap\ker\bar\partial^*.
\]
For the reverse inclusion, $C^*=\bar\partial^*\partial^*$, $E=\bar\partial\partial^*$ and $E^*=\partial\bar\partial^*$ all vanish under these three conditions. Elliptic theory on the compact manifold gives a finite-dimensional smooth kernel, its orthogonal projection $\mathcal H$ and the Green operator $\mathcal G$, with $\mathcal L\mathcal G=\operatorname{Id}-\mathcal H$.

Set $R=\mathcal L-C^*C$. Each term of $R$ has image in $\operatorname{im}\partial+\operatorname{im}\bar\partial$. Thus every smooth $(1,1)$-form $\alpha$ has the decomposition
\begin{equation}\label{eq:aeppli-projection-decomposition}
 \alpha=\mathcal H\alpha+R\mathcal G\alpha+C^*C\mathcal G\alpha.
\end{equation}
The first two terms lie in $\ker C$, since $C\partial=C\bar\partial=0$. The last is orthogonal to $\ker C$. Consequently
\begin{equation}\label{eq:aeppli-closed-projection}
 \mathcal P=\operatorname{Id}-C^*C\mathcal G
\end{equation}
is the orthogonal projection onto $\ker C$. If $C\alpha=0$, its last term in \eqref{eq:aeppli-projection-decomposition} vanishes, so its Aeppli class is represented by $\mathcal H\alpha$. A harmonic form in $\operatorname{im}\partial+\operatorname{im}\bar\partial$ is zero by orthogonality. This identifies $H_A^{1,1}$ with $\ker\mathcal L$ and proves finiteness.
\index{Aeppli Laplacian}

Now identify the varying type bundles and their $L^2$ spaces smoothly with those at $t=0$, using a unitary identification for the auxiliary metrics. The coefficients of $\mathcal L_t$ vary smoothly. Choose a circle $|z|=\varepsilon$ enclosing no nonzero eigenvalue of $\mathcal L_0$. Its resolvents are bounded from $H^s$ to $H^{s+4}$. The Neumann series for their perturbations gives smooth resolvents on this circle for small $t$, and the contour projections have constant rank. Constancy of $h_A^{1,1}$ makes that rank equal to $\dim\ker\mathcal L_t$. Hence the contour projections are precisely $\mathcal H_t$, with no additional small eigenvalues.

\Needspace{5\baselineskip}
The operators
\[
 \mathcal G_t=(\mathcal L_t+\mathcal H_t)^{-1}-\mathcal H_t
\]
vary smoothly on every Sobolev scale, by the same inverse argument and elliptic regularity. Formula \eqref{eq:aeppli-closed-projection} therefore gives the asserted smooth family $\mathcal P_t$. Complex conjugation sends $C_t$ to $-C_t$, exchanges $E_t$ with $E_t^*$ and exchanges the last two terms of \eqref{eq:aeppli-elliptic-operator}. Thus $\mathcal L_t$, $\mathcal H_t$, $\mathcal G_t$ and $\mathcal P_t$ commute with conjugation.
\end{proof}

\begin{theorem}\label{thm:pluriclosed-deformation-persistence}
Suppose that $(M,J_0)$ carries a pluriclosed Hermitian form $\omega_0$ and that $h_A^{1,1}(M,J_t)$ is constant near $0$. There is a smooth family of pluriclosed Hermitian forms $\omega_t$ on $(M,J_t)$ for small $t$, with the prescribed value $\omega_0$ at $t=0$.
\end{theorem}
\begin{proof}
Let $\pi_t^{1,1}$ denote type projection for $J_t$, and use the operators of \cref{lem:aeppli-smooth-projection}. Define
\[
 \omega_t=\mathcal P_t\bigl(\pi_t^{1,1}\omega_0\bigr).
\]
These are smooth real $(1,1)$-forms and satisfy $\partial_t\bar\partial_t\omega_t=0$. At $t=0$ the projection fixes $\omega_0$, since it already satisfies this equation. The tensors $\omega_t(\cdot,J_t\cdot)$ converge uniformly to the positive definite tensor $\omega_0(\cdot,J_0\cdot)$. Compactness gives a uniform positive lower bound at $t=0$, so they remain positive definite after shrinking $U$.
\end{proof}
\index{Pluriclosed metric!deformation}

The construction projects onto $\partial_t\bar\partial_t$-closed forms and obtains positivity by continuity near $\omega_0$. Closedness and curvature impose further equations. In particular, this persistence theorem supplies pluriclosed metrics along a given deformation; the balanced metric required by \cref{qu:balanced-pluriclosed-program} remains a separate geometric condition.
